% AI-written, AI cross-reviewed; no human sign-off; not a claim that RH is proved
% RH Notes blueprint source excerpts for N415
% Verbatim excerpts in the notation of the parent blueprint. The labelled editions and source ranges appear below.
% See the unit YAML for the current proof status. Earlier-round passages retain their original wording.
\section*{N415. Relative prime count in progressions}

% Parent source SHA-256: 44e3b079775c8d7cc0d43b0314bf43a3959be56b2bbafa7f26e5441a1a9163de
% blueprint/ch_T.tex lines 674 through 748; r5:lit:N415
\paragraph{A relative prime count and a character sum (N415, N420).}\label{r5:lit:N415}
Thorner and Zaman, \emph{Refinements to the prime number theorem for
arithmetic progressions} (Math. Z. 306 (2024), Article 54, \href{https://doi.org/10.1007/s00209-023-03414-3}{doi:10.1007/s00209-023-03414-3}), Corollary~1.4 and Remark~1.5, give the following
uniform formula for $q\ge2$, $(a,q)=1$ and $t\ge q^{12}$.
For the log-weighted prime count in the class $a$, the main term is
\[
 H_a(t)=\frac{t-s t^\beta/\beta}{\varphi(q)}
\]
and the error is $O(H_a(t)E_q(t))$, where $s=\chi_1(a)$ for the
exceptional real character and zero in its absence. For an absolute
$c>0$,
\[
 E_q(t)=e^{-c\log t/\log q}
       +\exp\{-c(\log t)^{3/5}/(\log\log t)^{1/5}\}.
\]
Their zero-free window at real height uses $1/(13\log(3q))$.
Equation~(1.2) and the discussion after it give
$1-\beta\gg q^{-1/2}$ when a zero exists.
See \href{https://arxiv.org/pdf/2108.10878\#page=3}{Thorner--Zaman,
Corollary 1.4 (p.~3), with (1.2) on p.~2}.

For N415 take $q=M$, a positive multiple of eight, and enlarge the
absolute constant $C_0$ in the hypothesis so that $X\ge q^{100}$.
On $[\sqrt X,X]$ partial summation applies to the formula at each $t$.
The main differential divided by $\log t$ is
$(1-s t^{\beta-1})\,dt/(\varphi(q)\log t)$.
Its integral gives the difference of the two logarithmic integrals in
\path{mainAP}, including its lower-end convention
$\mathrm{li}_2(x)=\int_2^x dt/\log t$.
Write $m(X)=1$ if $s\ne1$, and
$m(X)=\min(1,(1-\beta)\log X)$ if $s=1$. In the latter case integrate
over $[X^\beta,X]$ to see that $\mathrm{li}_2(X)-\mathrm{li}_2(X^\beta)
\asymp m(X)X/\log X$. In both cases $H_a(t)\ll m(X)t/\varphi(q)$ on
the stated interval, by $1-e^{-v}\le\min(1,v)$ in the cancelling case.
The endpoint error and its partial-summation integral are thus bounded by
$O(E_q(\sqrt X)\,m(X)X/(\varphi(q)\log X))$.
The initial prime count and lower-end main term are $O(\sqrt X)$.
Relative to the main term their cost is
$O(q^{3/2}\log X/\sqrt X)$ by (1.2). Both errors fit the two
exponentials of \path{PNTAPRelW} after reducing its absolute $c_3$.

Lean requires that a choice of no zero rules out every real
nonprincipal-character zero above $1-1/(8\log q)$. A zero selected
by the published window qualifies. If only the wider window contains a
candidate, take any such candidate. Then
$1-\beta\ge1/(13\log(3q))$, so the change in the main term, relative to
$\mathrm{li}_2(X)/\varphi(q)$, is
$O(e^{-c'\log X/\log q})$ for an absolute $c'>0$; by increasing $C_0$
it is smaller than half the original main term. The same relative
bound follows with the newly selected main term. This checks both
directions of the optional choice.
See \path{liX}, \path{piAP}, \path{IsExcPair}, \path{mainAP},
\path{PNTAPRelW} (\path{MeanSquareHypWIff.lean:24,26,29,33,38}) and
\path{pnt_progressions_relative_weak_holds} (\path{PNTFinal.lean:22}).

For N420 put $q=N\ge2$ and sum the log-weighted formula over
$a\in(\mathbb Z/q)^\times$ with multiplier $\chi(a)$. Orthogonality
centers the principal term and cancels the exceptional term unless
$\chi=\chi_1$. In the surviving case the hypothesis \path{NoExc}
forces $\beta\le 1-c_0$, so this term is $O(u^{1-c_0})$.
The sum of the absolute errors is $O(u E_q(u))$, since
$|s|\le1$ and the factor $\varphi(q)^{-1}$ cancels against the number
of classes. The additional prime powers contribute
$O(\sqrt u\,(\log u)^2)$. For $u\ge q^{12}$, $\lambda\ge\log q$
and $\lambda\ge1$, both terms are absorbed in
\[
 C u(\log u+\lambda)^2
   \exp\{-c\log u/(\lambda+\sqrt{\log u})\},
\]
with $c>0$ allowed to depend on $c_0$. For $e\le u<q^{12}$,
$\log u/(\lambda+\sqrt{\log u})\le12$ and the elementary
$O(u\log u)$ bound suffices. For the primitive principal character of
conductor one use modulus two, restore the omitted powers of two, and
use the trivial bound below $2^{12}$. This proves the requested
\path{QH.psi_thin} (\path{PsiThin.lean:385}).

